\chapter{Materials and Methods}
\label{ch:methods}

\section{Dataset}
The experiments use the \emph{arXiv co-authorship multiplex network}, a standard benchmark in which
nodes are authors and each of the 13 layers corresponds to one research domain (for example
\texttt{physics.soc-ph} or \texttt{cs.SI}). An edge in a layer means the two authors co-authored a
paper in that domain. The same author keeps the same node identity in every layer, which is exactly
the multiplex property~\cite{kivela2014} the model is designed to exploit. The 13 layers vary
considerably in size and density, which makes the network a good stress test for cross-layer methods.

\section{Preprocessing and Edge Splitting}
For each layer the observed edges are split into training, validation, and test sets. An equal number
of \emph{negative} edges (node pairs that are not connected) is sampled to balance each set, so that
link prediction becomes a balanced binary classification problem. The validation set is used only for
model selection and threshold tuning; the test set is untouched until final evaluation. The same
splits are reused across all methods so that every model is measured on identical data -- a
\textbf{fair-comparison} protocol.

\section{Node Feature Engineering}
Each node is described by six structural features computed from the graph -- among them its degree,
clustering coefficient, and centrality measures. These six numbers form the input vector for every
learning-based model, again ensuring that differences in performance come from the model and not from
the inputs.

\section{Baseline Methods}
Three families of baselines were implemented and evaluated under the shared protocol.
\begin{itemize}
  \item \textbf{Classical heuristics:} Common Neighbours, Jaccard, Adamic--Adar~\cite{adamic2003},
        and Preferential Attachment -- neighbourhood-based scores requiring no training.
  \item \textbf{Embedding baseline:} Node2Vec~\cite{node2vec}, which learns node vectors from biased
        random walks and predicts links from vector similarity.
  \item \textbf{Graph neural networks:} GCN~\cite{gcn} and GAT~\cite{gat}, trained per layer on the
        six node features, representing fixed-weight and attention-weighted message passing
        respectively.
\end{itemize}

\section{Proposed Model}
\label{sec:proposed}
The proposed model is an attention-based graph neural network built from four components that are
trained jointly. Figure~\ref{fig:pipeline} shows the complete pipeline, from the raw multiplex
network to the predicted link probability.

\begin{figure}[htbp]
\centering
\begin{tikzpicture}[
  node distance = 7mm and 12mm,
  box/.style   = {rectangle, rounded corners=2pt, draw=black, thick, align=center,
                  text width=7.2cm, minimum height=9mm, font=\small, fill=white},
  model/.style = {rectangle, rounded corners=2pt, draw=black, thick, align=center,
                  text width=7.2cm, minimum height=9mm, font=\small, fill=black!7},
  io/.style    = {rectangle, draw=black, very thick, align=center, text width=7.2cm,
                  minimum height=9mm, font=\small\bfseries, fill=black!14},
  arr/.style   = {-{Latex[length=2.4mm]}, thick}
]
\node[io]                 (data)  {arXiv Co-authorship Multiplex Network\\(authors as nodes, 13 research-domain layers)};
\node[box, below=of data] (split) {Edge splitting into train / validation / test\\ + balanced negative sampling};
\node[box, below=of split](feat)  {Node features (6 structural) \quad+\quad edge weak-tie signal\\ $s_{ij}=1-\mathrm{Jaccard}(i,j)$};
\node[model, below=of feat](enc)  {\textbf{(1) Weak-Tie-Aware GAT Encoder}\\ per-layer 64-dimensional node embeddings};
\node[model, below=of enc](cross) {\textbf{(2) Cross-Layer Attention}\\ learned, interpretable layer weights $\beta_\ell$};
\node[model, below=of cross](gate){\textbf{(3) Fusion Gate}\\ residual combine of within- and cross-layer evidence};
\node[model, below=of gate](dec)  {\textbf{(4) MLP Decoder}};
\node[io, below=of dec]   (out)   {Predicted link probability $\hat{y}_{ij}$};
\node[box, right=12mm of out, text width=3.3cm] (eval)
        {Evaluation:\\ ROC-AUC, PR-AUC, F1,\\ Hits@K, MRR\\ (5 seeds)};

\draw[arr] (data)  -- (split);
\draw[arr] (split) -- (feat);
\draw[arr] (feat)  -- (enc);
\draw[arr] (enc)   -- (cross);
\draw[arr] (cross) -- (gate);
\draw[arr] (gate)  -- (dec);
\draw[arr] (dec)   -- (out);
\draw[arr] (out)   -- (eval);

\begin{scope}[on background layer]
  \node[draw=black, dashed, rounded corners, fit=(enc)(cross)(gate)(dec),
        inner sep=3mm, label={[font=\small\itshape]above right:Proposed model}] {};
\end{scope}
\end{tikzpicture}
\caption{Overall methodology pipeline. Raw multiplex data is split and turned into node features and an
edge-level weak-tie signal; the four jointly trained components of the proposed model (dashed box)
produce a link probability, which is assessed with classification and ranking metrics over five seeds.}
\label{fig:pipeline}
\end{figure}

\subsection{Weak-Tie-Aware Attention Encoder}
The encoder is a graph-attention layer in the spirit of GAT~\cite{gat}, but with a crucial addition:
each edge carries a scalar \textbf{weak-tie signal} equal to $1-\text{Jaccard}$ of its two endpoints.
A high value marks a pair from different neighbourhoods -- a weak tie. By feeding this signal to the
attention mechanism (through the edge dimension of the attention layer), the encoder is told
\emph{where} the bridging links are, and can learn to represent them, instead of being blind to them
as neighbourhood heuristics are. The encoder outputs, for every node in every layer, a 64-dimensional
embedding derived from the six input features and the local weak-tie structure.

\subsection{Cross-Layer Attention}
The second component addresses the multiplex structure. For each node it collects the embeddings it
received in all 13 layers and combines them with a learned \textbf{attention weight}
$\beta_\ell$ per layer, following the cross-relation idea of HAN~\cite{han2019}. These weights are the
model's statement of how informative each research domain is, and they are read back as an
interpretable \emph{layer-importance} profile.

\subsection{Fusion Gate}
A gating step combines each node's own per-layer embedding with the cross-layer summary through a
learned residual gate, so that the model can rely on within-layer evidence, cross-layer evidence, or
a balance of the two, as the data dictates.

\subsection{Decoder}
Finally, a multilayer-perceptron decoder takes the fused embeddings of two candidate nodes and
outputs the probability that they are linked. An ablation variant, \emph{Proposed~B}, moves the
weak-tie signal into the decoder instead of the encoder, to isolate where the signal helps most.

\section{Evaluation Protocol}
Every method is evaluated per layer and averaged over the 13 layers. Classification quality is
measured by ROC-AUC, PR-AUC, and F1-score (with the threshold tuned on the validation set), and
ranking quality by Hits@K and MRR. To ensure the results are not an artefact of one random split,
each learning-based model is run with \textbf{five random seeds}; means and standard deviations are
reported, and a paired $t$-test assesses the significance of the improvement over the neural
baselines. A dedicated post-hoc analysis splits the test edges into \emph{weak-tie} and
\emph{common-neighbour} groups and reports ROC-AUC on each.

\section{Tools and Technologies}
The models were implemented in \textbf{Python} using \textbf{PyTorch} and \textbf{PyTorch~Geometric}
for the graph neural networks, \textbf{NetworkX} for graph construction and the classical heuristics,
\textbf{scikit-learn} for metrics and the Node2Vec pipeline, and \textbf{NumPy}/\textbf{pandas} for
data handling. Figures were produced with \textbf{Matplotlib}. Training used a single NVIDIA Quadro
P2000 GPU with early stopping on validation ROC-AUC.
